\cvsection{Selected Projects}
\begin{cventries}
  % \cventry
  %   {}
  %   {\vspace{-2em}Academic Projects}
  %   {}
  %   {September 2013 - May 2021}
  %   {
  %     \begin{cvitems}    
  %       \item {Applied {optimization} concepts to model and solve for optimal power-flow in a hypothetical smart-grid using {Julia} (see \href{https://github.com/bryanluu/smartgrid}{repository} on GitHub). The work was recognized as a \emph{Top Project} to be showcased to future students. \emph{(August 2019)}}
  %       % \item {Utilized Orbital Mechanics knowledge to simulate the Roche Limit in \textbf{Python} via Verlet Integration and \textbf{Object-Oriented Programming} (see \href{https://github.com/bryanluu/RocheLimit}{project repository} on GitHub). \emph{(March 2017)}}
  %       % \item Worked to design a laser upconversion setup for mid-IR to visible light for applications in Quantum Computing. Implemented a MATLAB script to perform spectroscopy using a monochrometer.
  %     \end{cvitems}
  %   }
  \cventry
    {Designer + Developer}
    {GRID Arcade}
    {Personal Project}
    {August 2025 - Present}
    {
      \begin{cvitems}
        \item Built a custom arcade in \textbf{C++} using an \textbf{Adafruit Metro M0} MCU  with an LED matrix display (see \href{https://github.com/bryanluu/GRID-Arcade}{GRID-Arcade} on GitHub). 
        \item Practiced \textbf{Object-Oriented} programming and \textbf{Memory-efficient Design} to implement different 4 different game-modes on limited hardware.
        \item Collaborated with a mechanical engineer to package system into a polished and robust product.
        \item Developed an emulation of software to debug control and game logic before hardware.
      \end{cvitems}
    }
  \cventry
    {\href{https://www.ubcsailbot.org/}{UBC Sailbot} (Student Team)}
    {Core Software Member \& Controls Lead}
    {UBC}
    {September 2013 -- August 2016}
    {
        \begin{cvitems}
            \item {Developed core \textbf{C++} control logic, \textbf{RS232} serial sensor forwarding, system simulation, along with power systems design.}
            \item Led research and development for new control system architectures, creating a boat simulation with components connected via \textbf{CAN-Bus}.
            \item Revamped the previous Thunderbird boat, improving the \textbf{Python} code structure and repairing electronics, ultimately
            resulting in a \emph{1st place, 50/50 finish at the 2014 International Robotic Sailing Regatta}.
            % \item Performed multiple presentations about the boat to industry professionals and potential sponsors, securing future investment and funding for the team.
        \end{cvitems}
    }
  % \iftoggle{resume}{}{\cventry
  %   {Co-project Manager \& Software Lead}
  %   {UBC ENPH 479 - Capstone Project}
  %   {Vancouver, BC, Canada}
  %   {Sept. 2017 - Present}
  %   {
  %     \begin{cvitems}
  %       \item Worked under Dr. Jeff Young and Dr. Jingda Wu to design a laser upconversion setup for mid-IR to visible light for applications in Quantum Computing.
  %       \item Implemented a MATLAB script to perform spectroscopy using a monochrometer.
  %       \item Assisted in choosing, ordering, and aligning optical components for the setup.
  %     \end{cvitems}
  %   }}
  % \cventry
  %   {Core Software Member \& Controls Lead in 2016}
  %   {\href{https://www.ubcsailbot.org/}{UBC Sailbot}}
  %   {Vancouver, BC, Canada}
  %   {Sept. 2013 - Aug. 2016}
  %   {
  %     \begin{cvitems}
  %       \item Employed \textbf{Agile} and \textbf{Test-Driven} development practices to implement core sailing logic in \textbf{C++}, including sensor data forwarding over \textbf{RS232} for transatlantic boat \href{https://www.ubcsailbot.org/new-page-2}{Ada}.
  %       \item Revamped the previous Thunderbird boat, improving the \textbf{Python} code structure and repairing electronics, ultimately
  %       resulting in a \emph{1st place, 50/50 finish at the 2014 International Robotic Sailing Regatta}.
  %       \item Performed multiple presentations about the boat to industry professionals and potential sponsors, securing future investment and funding for the team.
  %     \end{cvitems}
  %   }
  % \iftoggle{resume}{}{\cventry
  %   {Attendee}
  %   {USEQIP (Undergraduate School for Experimental Quantum Information Processing)}
  %   {Waterloo, ON, Canada}
  %   {June 2015}
  %   {
  %     \begin{cvitems}
  %       \item Selected among 100+ undergraduates to attend an exclusive 2-week program to learn about experimental Quantum Information.
  %       \item Learned basic quantum algorithms including Shor's Algorithm, Deutsch-Jozsa Algorithm, and the Quantum Fourier Transform.
  %       \item Successfully applied the Quantum Fourier Transform onto a two-qubit NMR system.
  %       \item Received training on the basics of using a nanofabrication cleanroom facility.
  %     \end{cvitems}
  %   }}
%   \cventry
%     {Software Lead}
%     {UBC ENPH - Summer Robotics Competition}
%     {Vancouver, BC, Canada}
%     {May 2014 - Aug. 2014}
%     {
%       \begin{cvitems}
%         \item Worked in a team of 4 to design and prototype a fully autonomous robot for the annual Engineering Physics
%         Robotics Competition, culminating in a \emph{3rd place finish out of 56 students, 14 teams}.
%         \item Coordinated software development as Software Lead, designing the software architecture with \textbf{C++} on an \textbf{Atmel ATmega128} based microprocessor.
%         \item Programmed extensive unit tests in \textbf{CppUnit} to validate external \textbf{PID control} systems used for steering.
%         \item Assembled and designed major electrical circuits, such as the IR detection and homing circuit as well as the Motherboard, responsible for delivering power to all main electronics.
%         % \item Devised robot’s gear systems, using SolidWorks CAD Software and a laser cutter for prototyping.
%         \item Managed \href{https://github.com/bryanluu/Fizz-253}{team repository} on GitHub, and produced the \href{http://indyontherocks.wixsite.com/enph253}{website} and \href{https://youtu.be/znlS3rSShug}{video} documenting project work; these were demonstrated to incoming first year students.
%       \end{cvitems}
%     }
%   \cventry
%     {Student}
%     {Software Modeling}
%     {Vancouver, BC, Canada}
%     {Sept. 2008 - May 2018}
%     {
%       \begin{cvitems}
%         \item Implemented various numerical integration schemes (Trapezium,Gauss,etc.) for my Numerical Analysis class, including modelling of the mold-filling problem.
%         \item Combined MATLAB and Arduino programming with Finite Differences Approximation to model and analyze the effects of heat transfer through an Aluminum Rod for Thermodynamics Lab course.
%         \item Expanded MATLAB GUI skills to create an application for visualizing parametric curves in 3D space, as an extracurricular aid to my Vector Calculus course. Project found at \href{https://github.com/Fizz-2013/MATLAB-Projects}{https://github.com/Fizz-2013/MATLAB-Projects}.
%         \item In Grade 9, attempted to create 3D graphics with 2D Java engine by manipulating lines to create single-focus perspective, combining both passions in visual art and computer science.
%       \end{cvitems}
%     }
\end{cventries}
